\section{The set of rational numbers}
\label{sec:rational-numbers}

The exercises in this section ask whether particular numbers belong to
$\Q$. Recall that a number is rational when it can be written as $m/n$,
where $m,n\in\Z$ and $n\ne0$. A number is irrational when it is a real
number that is not rational. The following elementary tools explain the
steps used in the solutions.

\step{Tool 1: Why rational arithmetic stays rational}
Suppose $r=m/n$ and $s=p/q$, where $m,n,p,q\in\Z$ and $n,q\ne0$.
Putting fractions over a common denominator gives
\[
 r+s=\frac{mq+pn}{nq},\qquad
 r-s=\frac{mq-pn}{nq},\qquad
 rs=\frac{mp}{nq}.
\]
Each numerator is an integer, and $nq$ is a nonzero integer. Each displayed
quantity is therefore rational by the definition of $\Q$. If $s\ne0$,
then $p\ne0$, and
\[
 \frac{r}{s}=\frac{mq}{np}
\]
also has an integer numerator and a nonzero integer denominator. Thus
division by a nonzero rational number preserves rationality. Integers
are rational because $j=j/1$. In particular, multiplication by an integer
and subtraction of an integer preserve rationality.

For a positive integer $k$, repeated multiplication gives
\[
 r^k=\underbrace{r\cdot r\cdots r}_{k\text{ factors}}
     =\frac{m^k}{n^k}\in\Q,
\]
since $m^k,n^k\in\Z$ and $n^k\ne0$. These formulas justify every use of
``rational closure'' below.

\step{Tool 2: The integer cancellation fact behind the rational zeros theorem}
Two integers are relatively prime when they have no common positive
divisor other than $1$. We use this elementary fact: if $u,v\in\Z$ are
relatively prime, $v\ne0$, and $v$ divides $uw$, then $v$ divides $w$.
Here ``$v$ divides $z$'' means that $z=vk$ for some $k\in\Z$.

For completeness, the cancellation fact follows from a short integer
argument. The positive integers of the form $Au+Bv$, with $A,B\in\Z$,
form a nonempty set: one of $v$ and $-v$ is positive. Let $d$ be the
least such positive integer, and write $d=A_0u+B_0v$. The division
algorithm gives $u=hd+t$, with $h\in\Z$ and $0\le t<d$. Then
\[
 t=u-hd=(1-hA_0)u-hB_0v.
\]
If $t>0$, this would be a smaller positive integer of the same form,
contradicting the choice of $d$. Hence $t=0$, so $d$ divides $u$.
Applying the same argument to division of $v$ by $d$ shows that $d$
divides $v$. Relative primality forces $d=1$. Consequently there are
integers $A_0,B_0$ with
\[
 A_0u+B_0v=1.
\]
Multiplying this identity by $w$ gives
\[
 w=A_0(uw)+B_0vw.
\]
If $v$ divides $uw$, both terms on the right are integer multiples of
$v$, so $w$ is also an integer multiple of $v$. This proves the
cancellation fact. The integer division algorithm and the existence
of a least positive integer are the only integer facts used here.

One consequence needed below is worth spelling out. If $u,v$ are
relatively prime and $v$ divides $u^k$, then cancellation in
$u^k=u\,u^{k-1}$ shows that $v$ divides $u^{k-1}$. Repeating this step
eventually shows that $v$ divides $1$. When $u\ne0$, the last power in
this argument is $u^0=1$. When $u=0$, relative primality already forces
$\abs{v}=1$, so the same conclusion holds directly.

\step{Tool 3: A monic polynomial can have only integer rational zeros}
A polynomial is \emph{monic} when the coefficient of its highest power
is $1$. Suppose
\[
 P(x)=x^d+a_{d-1}x^{d-1}+\cdots+a_1x+a_0,
 \qquad a_0,\ldots,a_{d-1}\in\Z,
\]
where $d\ge1$. Assume that $r\in\Q$ and $P(r)=0$. Write
$r=m/n$ in lowest terms, with $n>0$. Thus $m,n\in\Z$ are relatively
prime. If $m=0$, then $r=0$ is already an integer. For the denominator
argument we may therefore assume $m\ne0$. Substituting and multiplying
by $n^d$ gives
\[
 m^d+\sum_{j=0}^{d-1}a_jm^j n^{d-j}=0.
\]
Every term in the sum contains a factor $n$. Equivalently,
\[
 m^d=-n\left(\sum_{j=0}^{d-1}a_jm^j n^{d-j-1}\right).
\]
The expression in parentheses is an integer, so $n$ divides $m^d$.
Tool~2, applied repeatedly to the relatively prime integers $m,n$,
shows that $n$ divides $1$. As $n$ is positive, $n=1$. Therefore
$r=m\in\Z$.

There is a further restriction if $a_0\ne0$. In that case $r\ne0$,
because $P(0)=a_0\ne0$. Substituting the integer $m$ into $P(m)=0$
and collecting all terms other than the constant gives
\[
 a_0=-m\left(m^{d-1}+\sum_{j=1}^{d-1}a_jm^{j-1}\right).
\]
For $d=1$ the sum is empty and equals $0$. In every case the expression
in parentheses is an integer. Hence $m$ divides $a_0$.

We have proved the monic case of Ross's Rational Zeros Theorem
(Theorem~2.2 and Corollary~2.3):
\begin{quote}
A rational zero of a monic polynomial with integer coefficients must
be an integer. If the constant coefficient is nonzero, that integer
must divide the constant coefficient.
\end{quote}
In particular, if $N$ is a positive integer and $k$ is a positive integer,
then a rational solution of $x^k=N$ must be an integer, because it is
a zero of the monic integer polynomial $x^k-N$.

\step{Tool 4: Bounding a positive root by consecutive integers}
For $0\le u<v$ and a positive integer $k$,
\[
 v^k-u^k=(v-u)
 \bigl(v^{k-1}+v^{k-2}u+\cdots+u^{k-1}\bigr)>0.
\]
For $k=1$ the factor in parentheses is $1$. For $k\ge2$, all its
terms are nonnegative and its first term $v^{k-1}$ is positive.
Thus raising nonnegative numbers to the $k$th power preserves their
strict order. This is an algebraic fact; no limiting argument is needed.

Let $x$ be the nonnegative $k$th root of a positive integer $N$.
Suppose a nonnegative integer $a$ satisfies
\[
 a^k<N<(a+1)^k.
\]
If $x\le a$, then $x^k\le a^k<N$, contradicting $x^k=N$.
If $x\ge a+1$, then $x^k\ge(a+1)^k>N$, giving the same contradiction.
It follows that
\[
 a<x<a+1.
\]
There is no integer strictly between the consecutive integers $a$ and
$a+1$. Tool~3 therefore implies that $x$ is irrational. This gives the
criterion used in the first two exercises:
\begin{quote}
If a positive integer lies strictly between two consecutive nonnegative
integer $k$th powers, its positive $k$th root is irrational.
\end{quote}

\exercise{2.1}{Five square roots of integers}
\goal Prove that $\sqrt3$, $\sqrt5$, $\sqrt7$, $\sqrt{24}$, and
$\sqrt{31}$ are irrational.

\idea Each number is a root of a monic integer polynomial $x^2-N$.
If it were rational, Tool~3 would make it an integer. We locate each
root strictly between consecutive integers to rule that out.

\step{1. Record the exact integer-square comparisons}
The needed computations are
\[
 1^2=1,\quad 2^2=4,\quad 3^2=9,\quad
 4^2=16,\quad 5^2=25,\quad 6^2=36.
\]
They give the following bounds by Tool~4:
\begin{center}
\begin{tabular}{@{}ccc@{}}
\toprule
$N$ & Consecutive square bounds & Consequent root bounds\\
\midrule
$3$  & $1^2=1<3<4=2^2$   & $1<\sqrt3<2$\\
$5$  & $2^2=4<5<9=3^2$   & $2<\sqrt5<3$\\
$7$  & $2^2=4<7<9=3^2$   & $2<\sqrt7<3$\\
$24$ & $4^2=16<24<25=5^2$ & $4<\sqrt{24}<5$\\
$31$ & $5^2=25<31<36=6^2$ & $5<\sqrt{31}<6$\\
\bottomrule
\end{tabular}
\end{center}

\step{2. Apply the rational-root argument to each number}
\begin{enumerate}[label=(\alph*)]
\item Suppose $\sqrt3$ were rational. Since
$(\sqrt3)^2-3=0$, Tool~3 applied to $x^2-3$ would imply
$\sqrt3\in\Z$. But $1<\sqrt3<2$, and no integer satisfies those
two inequalities. This contradiction proves $\sqrt3\notin\Q$.

\item Suppose $\sqrt5$ were rational. The equation
$(\sqrt5)^2-5=0$ and Tool~3 imply $\sqrt5\in\Z$. The bounds
$2<\sqrt5<3$ exclude every integer, so $\sqrt5\notin\Q$.

\item Suppose $\sqrt7$ were rational. The monic integer equation
$(\sqrt7)^2-7=0$ would make $\sqrt7$ an integer. The strict bounds
$2<\sqrt7<3$ make this impossible. Hence $\sqrt7\notin\Q$.

\item Suppose $\sqrt{24}$ were rational. Because it satisfies
$x^2-24=0$, it would be an integer by Tool~3. Its bounds
$4<\sqrt{24}<5$ contradict that conclusion. Thus
$\sqrt{24}\notin\Q$.

\item Suppose $\sqrt{31}$ were rational. Tool~3 applied to
$x^2-31$ would make it an integer. Since $5<\sqrt{31}<6$, there is
no such integer. Therefore $\sqrt{31}\notin\Q$.
\end{enumerate}

\step{3. Establish one auxiliary case for later exercises}
We will also need that $\sqrt2$ is irrational. If it were rational,
the equation $(\sqrt2)^2-2=0$ would make it an integer. But
\[
 1^2=1<2<4=2^2
 \quad\Longrightarrow\quad 1<\sqrt2<2,
\]
which is impossible for an integer. Thus $\sqrt2\notin\Q$ as well.

\answer{All five requested square roots are irrational. The same argument
also proves the auxiliary fact $\sqrt2\notin\Q$.}

\exercise{2.2}{Three higher roots of integers}
\goal Prove that $\sqrt[3]{2}$, $\sqrt[7]{5}$, and $\sqrt[4]{13}$
are irrational.

\idea Tool~3 applies to $x^3-2$, $x^7-5$, and $x^4-13$.
We compute suitable consecutive integer powers, then apply Tool~4.

\step{1. The cube root of $2$}
Let $x=\sqrt[3]{2}$. By the definition of the positive cube root,
$x>0$ and $x^3=2$. Assume that $x\in\Q$. Since $x$ is a zero of
the monic polynomial $x^3-2$, whose coefficients are integers,
Tool~3 gives $x\in\Z$.

Now $1^3=1$ and $2^3=2\cdot2\cdot2=8$. Consequently
\[
 1^3=1<2<8=2^3.
\]
Tool~4 gives $1<x<2$. No integer lies in that interval, contradicting
$x\in\Z$. Hence $\sqrt[3]{2}$ is irrational.

\step{2. The seventh root of $5$}
Let $y=\sqrt[7]{5}$. Then $y>0$ and $y^7=5$. If $y$ were rational,
Tool~3 applied to $x^7-5$ would imply $y\in\Z$. The exact powers are
\[
 1^7=1,\qquad
 2^7=2\cdot2\cdot2\cdot2\cdot2\cdot2\cdot2=128.
\]
Thus
\[
 1^7=1<5<128=2^7,
\]
so $1<y<2$. This excludes all integers and contradicts the consequence
of the rationality assumption. Therefore $\sqrt[7]{5}$ is irrational.

\step{3. The fourth root of $13$}
Let $z=\sqrt[4]{13}$, where the fourth-root notation denotes the
positive real root. Then $z>0$ and $z^4=13$. If $z\in\Q$, the monic
integer equation $z^4-13=0$ gives $z\in\Z$ by Tool~3. However,
\[
 1^4=1<13<16=2\cdot2\cdot2\cdot2=2^4.
\]
It follows that $1<z<2$, contradicting integer membership. Hence
$\sqrt[4]{13}$ is irrational.

\answer{The numbers $\sqrt[3]{2}$, $\sqrt[7]{5}$, and
$\sqrt[4]{13}$ are all irrational.}

\exercise{2.3}{A nested square root}
\goal Prove that $\sqrt{2+\sqrt2}$ is irrational.

\idea Assume the outer square root is rational. Square it and isolate
$\sqrt2$. Rational closure would then force $\sqrt2$ to be rational,
contradicting the auxiliary case in Exercise~2.1.

\step{1. Name the number and record its defining equation}
Put
\[
 a=\sqrt{2+\sqrt2}.
\]
The inner root $\sqrt2$ is positive, so $2+\sqrt2>0$. The outer square
root therefore denotes a positive real number. Its defining equation is
\[
 a^2=2+\sqrt2.
\]

\step{2. Assume rationality and rearrange}
Suppose, for a contradiction, that $a\in\Q$. Then $a^2\in\Q$ by
rational closure under multiplication. Subtracting the rational integer
$2$ from both sides gives
\[
 \sqrt2=a^2-2\in\Q.
\]
In fraction form, if $a=m/n$ with $m,n\in\Z$ and $n\ne0$, this is
\[
 a^2-2=\frac{m^2}{n^2}-\frac{2n^2}{n^2}
       =\frac{m^2-2n^2}{n^2},
\]
an integer divided by a nonzero integer. Thus the rationality conclusion
is justified directly by the definition, as well as by Tool~1.

\step{3. Identify the contradiction}
Exercise~2.1 established $\sqrt2\notin\Q$. The rationality assumption
for $a$ has produced its negation, $\sqrt2\in\Q$. Therefore the
assumption is false.

\answer{$\sqrt{2+\sqrt2}$ is irrational.}

\exercise{2.4}{A nested cube root}
\goal Prove that $\sqrt[3]{5-\sqrt3}$ is irrational.

\idea Cube the number to remove the outer radical, then isolate
$\sqrt3$. If the original number were rational, this rearrangement
would express $\sqrt3$ using rational operations only.

\step{1. Check the sign and write the cube equation}
Exercise~2.1 gives $1<\sqrt3<2$. In particular,
\[
 5-\sqrt3>5-2=3>0.
\]
Thus
\[
 a=\sqrt[3]{5-\sqrt3}
\]
is a positive real number, and the definition of cube root gives
\[
 a^3=5-\sqrt3.
\]

\step{2. Derive a rational expression for $\sqrt3$ under the assumption}
Assume $a\in\Q$. Then $a^3\in\Q$ by Tool~1. Adding $\sqrt3$ to the
cube equation and subtracting $a^3$ gives
\[
 a^3+\sqrt3=5,
 \qquad\text{hence}\qquad
 \sqrt3=5-a^3.
\]
Both $5$ and $a^3$ are rational, so their difference is rational.
Explicitly, if $a=m/n$ with integer $m,n$ and $n\ne0$, then
\[
 5-a^3=\frac{5n^3-m^3}{n^3}\in\Q.
\]

\step{3. Finish the contradiction}
The preceding step gives $\sqrt3\in\Q$, contrary to Exercise~2.1.
Therefore $a\notin\Q$.

\answer{$\sqrt[3]{5-\sqrt3}$ is irrational.}

\exercise{2.5}{A positive real fractional power}
\goal Prove that $(3+\sqrt2)^{2/3}$ is irrational.

\idea First explain the fractional exponent using the positive real
cube root. Cubing the proposed number then produces a simple expression
$11+6\sqrt2$, from which $\sqrt2$ can be isolated.

\step{1. Interpret the exponent without a sign ambiguity}
Since $\sqrt2>0$, the base
\[
 u=3+\sqrt2
\]
is positive. Let $t=\sqrt[3]{u}>0$. Thus $t^3=u$, and the positive real
fractional power in the problem is
\[
 a=u^{2/3}=(\sqrt[3]{u})^2=t^2.
\]
It follows from integer power rules that
\[
 a^3=(t^2)^3=t^6=(t^3)^2=u^2=(3+\sqrt2)^2.
\]
This explains precisely why cubing removes the exponent $2/3$.

\step{2. Expand the square completely}
Using $(v+w)^2=v^2+2vw+w^2$ gives
\begin{align*}
 a^3
  &=(3+\sqrt2)^2\\
  &=3^2+2\cdot3\cdot\sqrt2+(\sqrt2)^2\\
  &=9+6\sqrt2+2\\
  &=11+6\sqrt2.
\end{align*}
Subtracting $11$ and dividing by $6$ therefore yields
\[
 a^3-11=6\sqrt2,
 \qquad
 \sqrt2=\frac{a^3-11}{6}.
\]

\step{3. Use rational closure to obtain the contradiction}
Suppose $a\in\Q$. Then $a^3\in\Q$, hence $a^3-11\in\Q$.
Because $6$ is a nonzero rational integer, division by $6$ preserves
rationality. The last displayed equation would imply $\sqrt2\in\Q$.
This contradicts the auxiliary irrationality result in Exercise~2.1.

\answer{$(3+\sqrt2)^{2/3}$, with the positive real interpretation of the
fractional power, is irrational.}

\exercise{2.6}{A rational quadratic expression}
\goal Explain why $b\in\Q$ implies $4-7b^2\in\Q$.

\idea Start with the defining fraction for a rational number and write
$4-7b^2$ as one fraction. Check its numerator and denominator separately.

\step{1. Write the rationality hypothesis explicitly}
Assume $b\in\Q$. By definition, there are integers $m,n$ such that
\[
 b=\frac mn,\qquad n\ne0.
\]
The numerator $m$ may be positive, negative, or zero. No restriction on
the sign of $b$ is needed for this argument.

\step{2. Perform the arithmetic over one denominator}
Squaring and multiplying gives
\[
 b^2=\frac{m^2}{n^2},
 \qquad 7b^2=\frac{7m^2}{n^2}.
\]
As $n\ne0$, its square $n^2$ is nonzero, so we may also write
$4=4n^2/n^2$. Subtracting now gives
\[
 4-7b^2
 =\frac{4n^2}{n^2}-\frac{7m^2}{n^2}
 =\frac{4n^2-7m^2}{n^2}.
\]
The integer operations of squaring, multiplication, and subtraction
show that $4n^2-7m^2\in\Z$. The denominator $n^2$ is a nonzero
integer. The displayed fraction is therefore rational by definition.

\step{3. Explain the connection with Example 6}
The same reasoning can be summarized using Tool~1: $b\in\Q$ implies
$b^2\in\Q$, then $7b^2\in\Q$, and then $4-7b^2\in\Q$.
This is the rational-closure step used in Example~6. For the particular
number considered there, the radical equation also gives
\[
 7b^2=4-2\sqrt3,
 \qquad\text{so}\qquad
 4-7b^2=2\sqrt3.
\]
If that $b$ were rational, the result just proved would make
$2\sqrt3$ rational. Dividing by the nonzero rational number $2$ would
make $\sqrt3$ rational, contradicting Exercise~2.1. This shows how the
closure statement supplies the needed step in an irrationality proof.

\answer{For every rational $b=m/n$ with $n\ne0$,
$4-7b^2=(4n^2-7m^2)/n^2\in\Q$.}

\exercise{2.7}{Two radicals that simplify to rational numbers}
\goal Show that
\[
 \text{(a)}\quad \sqrt{4+2\sqrt3}-\sqrt3,
 \qquad
 \text{(b)}\quad \sqrt{6+4\sqrt2}-\sqrt2
\]
are rational, and find their values. In each expression the final
subtracted square root lies outside the first radical.

\idea Recognize each radicand as the square of a positive sum. Positivity
is essential: for a real $v$, the principal square root satisfies
$\sqrt{v^2}=\abs{v}$. To replace that absolute value by $v$, we check
that $v>0$.

\step{1. Part (a): Expand the proposed positive square root}
Consider $v=1+\sqrt3$. Because $\sqrt3>0$, we have $v>0$. Expanding
its square gives
\begin{align*}
 (1+\sqrt3)^2
  &=1^2+2\cdot1\cdot\sqrt3+(\sqrt3)^2\\
  &=1+2\sqrt3+3\\
  &=4+2\sqrt3.
\end{align*}
The radicand is positive, and $1+\sqrt3$ is a positive number whose
square equals that radicand. It is therefore the principal square root:
\[
 \sqrt{4+2\sqrt3}=1+\sqrt3.
\]
One can also justify the choice of sign from uniqueness. If two
nonnegative numbers $v,w$ have the same positive square, then
$(v-w)(v+w)=v^2-w^2=0$. Since $v+w>0$, it follows that $v-w=0$,
so $v=w$. There is thus only one nonnegative square root here.

Substituting the identity into the expression in part (a) gives
\[
 \sqrt{4+2\sqrt3}-\sqrt3
 =(1+\sqrt3)-\sqrt3
 =1+(\sqrt3-\sqrt3)
 =1+0=1.
\]
As $1=1/1$, this value is rational.

\step{2. Part (b): Expand the second proposed positive square root}
Now consider $w=2+\sqrt2$. It is positive because $\sqrt2>0$.
Its square expands as
\begin{align*}
 (2+\sqrt2)^2
  &=2^2+2\cdot2\cdot\sqrt2+(\sqrt2)^2\\
  &=4+4\sqrt2+2\\
  &=6+4\sqrt2.
\end{align*}
Thus $2+\sqrt2$ is the positive square root of $6+4\sqrt2$.
The sign check permits the identity
\[
 \sqrt{6+4\sqrt2}=2+\sqrt2.
\]
It follows that
\[
 \sqrt{6+4\sqrt2}-\sqrt2
 =(2+\sqrt2)-\sqrt2
 =2+(\sqrt2-\sqrt2)
 =2+0=2.
\]
Since $2=2/1$, this value is rational as well.

\answer{Part (a) equals $1$ and part (b) equals $2$; both expressions
are rational numbers.}

\exercise{2.8}{All rational zeros of an eighth-degree polynomial}
\goal Find every rational solution of
\[
 x^8-4x^5+13x^3-7x+1=0.
\]

\idea The polynomial is monic, has integer coefficients, and has constant
term $1$. The monic rational-zero result restricts every rational root
to an integer divisor of $1$. We must test both candidates and explain
why this exhausts the rational possibilities.

\step{1. Obtain the complete candidate list}
Let
\[
 P(x)=x^8-4x^5+13x^3-7x+1.
\]
The coefficient of $x^8$ is $1$, so $P$ is monic. Its missing powers
have coefficient $0$, which is also an integer. Its constant coefficient
is $1\ne0$. Thus Tool~3 applies to every rational zero of $P$.
It says that such a zero must be an integer dividing $1$. If an integer
$m$ divides $1$, then $1=mk$ for some integer $k$. Taking absolute
values gives $1=\abs{m}\abs{k}$, so both nonnegative integer factors
must equal $1$. In particular $\abs{m}=1$, and hence $m=1$ or $m=-1$.
The full rational candidate list is therefore
\[
 \{1,-1\}.
\]

Here is the denominator step directly for this particular polynomial.
If a rational root is $m/n$ in lowest terms, with $n>0$, multiplication
of $P(m/n)=0$ by $n^8$ gives
\[
 m^8-4m^5n^3+13m^3n^5-7mn^7+n^8=0.
\]
Solving for $m^8$ gives
\[
 m^8=n\bigl(4m^5n^2-13m^3n^4+7mn^6-n^7\bigr),
\]
so $n$ divides $m^8$. Since $m,n$ are relatively prime, repeated
coprime cancellation gives $n\mid1$, hence $n=1$.
After this reduction, $P(m)=0$ can be rearranged as
\[
 1=-m^8+4m^5-13m^3+7m
   =m\bigl(-m^7+4m^4-13m^2+7\bigr).
\]
This proves $m\mid1$ and recovers precisely the two candidates above.
Also $P(0)=1$, so zero has not been omitted as a possible root.

\step{2. Evaluate the candidate $1$}
Every positive power of $1$ equals $1$. Consequently
\begin{align*}
 P(1)
 &=1^8-4\cdot1^5+13\cdot1^3-7\cdot1+1\\
 &=1-4+13-7+1\\
 &=-3+13-7+1\\
 &=10-7+1\\
 &=3+1=4.
\end{align*}
Since $4\ne0$, the number $1$ is not a solution.

\step{3. Evaluate the candidate $-1$}
The eighth power of $-1$ is $1$, while its fifth and third powers
are $-1$. Substituting with the signs retained gives
\begin{align*}
 P(-1)
 &=(-1)^8-4(-1)^5+13(-1)^3-7(-1)+1\\
 &=1-4(-1)+13(-1)+7+1\\
 &=1+4-13+7+1\\
 &=5-13+7+1\\
 &=-8+7+1\\
 &=-1+1=0.
\end{align*}
Thus $-1$ is a solution, and it is rational because $-1=(-1)/1$.

\step{4. Explain why the search is complete}
Every rational solution must belong to the candidate set $\{1,-1\}$
from Step~1. Step~2 rejects $1$, and Step~3 verifies $-1$. There are
no other rational candidates to examine.

\answer{The set of all rational solutions is $\{-1\}$.}
